Este capítulo detalha uma arquitetura orientada a \textit{Complex Event Processing} (CEP) para captura, integração e organização de eventos que alimentam o gêmeo digital em contextos PBL. A modelagem alinha-se às visões estruturais, comportamentais e de processo definidas pela ISO/IEC/IEEE 42010 \cite{iso42010} e aos viewpoints do RM‑ODP \cite{iso10746}, com ênfase em observabilidade, rastreabilidade e qualidades arquiteturais (latência, confiabilidade, interoperabilidade e portabilidade), sem qualquer mensuração de desempenho discente.

Os objetivos deste capítulo são:
\begin{itemize}
  \item Definir um esquema de eventos e respectivos contratos de dados;
  \item Estabelecer chaves e encadeamentos para rastreabilidade fim a fim;
  \item Descrever a arquitetura de ingestão/CEP e táticas arquiteturais;
  \item Derivar vistas e sinais operacionais para apoio à decisão docente.
\end{itemize}

\subsection{Princípios de CEP e Escopo}

O CEP estrutura o processamento de acontecimentos discretos (eventos) em janelas temporais, correlacionando múltiplas fontes para identificar padrões e produzir resultados derivados (sinais) de forma tempestiva. Nesta proposta:
\begin{itemize}
  \item \textbf{Imutabilidade}: eventos são registros imutáveis do que ocorreu;
  \item \textbf{Semântica temporal}: prioriza-se \textit{event time}; \textit{processing time} é registrado para auditoria;
  \item \textbf{Ordenação e \textit{watermarks}}: reordenação limitada; eventos atrasados são tratados com políticas explícitas;
  \item \textbf{Idempotência}: escrita e reprocessamento idempotentes em todas as etapas;
  \item \textbf{Compatibilidade de esquemas}: contratos versionados com políticas de \textit{backward/forward} compatibility;
  \item \textbf{Privacidade e minimização}: separação de PII, classificação de confidencialidade e acesso por papéis.
\end{itemize}

\subsection{Caso 1 — Esquema de Eventos e Contratos}

\textbf{Taxonomia de eventos} (exemplos):
\begin{itemize}
  \item \textit{Técnicos}: commits, pull requests, \textit{builds}, testes, releases;
  \item \textit{Operacionais}: plannings, dailies, reviews, retrospectivas;
  \item \textit{Acadêmicos}: aulas, presença, APs (registro de ocorrência, não de nota);
  \item \textit{Documentais}: TAPI, requisitos, artefatos entregues.
\end{itemize}

\textbf{Envelope mínimo} (campos ilustrativos):
\begin{itemize}
  \item \texttt{event\_id} (UUID), \texttt{trace\_id}, \texttt{event\_type}, \texttt{schema\_version};
  \item \texttt{event\_time}, \texttt{ingested\_at} (ISO 8601);
  \item \texttt{source} (ex.: git, lms, tracker), \texttt{source\_id};
  \item \texttt{actor\_id}, \texttt{team\_id}, \texttt{module\_id}, \texttt{sprint\_id};
  \item \texttt{pii\_tags}, \texttt{confidentiality}, \texttt{checksum}.
\end{itemize}

\textbf{Contratos de dados}:
\begin{itemize}
  \item Definição por JSON Schema com registro de versões e \textit{lifecycle} de depreciação;
  \item Validação na borda (conectores) e no \textit{landing} (bronze), com relatórios de rejeição;
  \item Estratégias de compatibilidade: campos opcionais, mapas de transformação e preenchimentos padrão.
\end{itemize}

\textbf{Camadas de dados}:
\begin{itemize}
  \item \textit{Bronze}: eventos crus, particionados por data/fonte;
  \item \textit{Silver}: normalização por contratos, \textit{dedup}, ordenação por \textit{event time};
  \item \textit{Gold}: materializações de leitura e \textit{denormalizações} por sinais/relatórios.
\end{itemize}

\subsection{Caso 2 — Rastreabilidade Fim a Fim e Identidades}

\textbf{Espaço de chaves}: \texttt{person\_id}, \texttt{team\_id}, \texttt{module\_id}, \texttt{sprint\_id}, \texttt{artifact\_id}, \texttt{requirement\_id}, \texttt{repo\_id}, \texttt{pr\_id}, \texttt{commit\_sha}.

\textbf{Encadeamentos}:
\begin{itemize}
  \item Do evento ao recurso: \texttt{event\_id} → (\texttt{commit\_sha}|\texttt{pr\_id}|\texttt{artifact\_id});
  \item Do indivíduo ao contexto: \texttt{person\_id} → \texttt{team\_id} → \texttt{module\_id} → \texttt{sprint\_id};
  \item Da intenção ao produto: \texttt{requirement\_id} → backlog → PR/commit → artefato.
\end{itemize}

\textbf{Governança e PII}:
\begin{itemize}
  \item Tabelas de correspondência segregadas (reversíveis apenas por papéis autorizados);
  \item Minimização por \textit{hash}/pseudonimização quando aplicável;
  \item Auditoria e linhagem: trilhas de transformação, \textit{data provenance} e \textit{lineage} consultáveis.
\end{itemize}

\subsection{Caso 3 — Arquitetura de Ingestão/CEP e Táticas Arquiteturais}

\textbf{Componentes}: conectores de fontes (Git, tracker, LMS), fila/mensageria, orquestrador de ingestão, motor CEP (janelas \textit{tumbling/sliding/session}, \textit{joins} por chaves), camadas bronze/silver/gold, APIs de leitura e materializações.

\textbf{Táticas arquiteturais} (exemplos):
\begin{itemize}
  \item \textit{Backpressure} e controle de taxa; \textit{retry} exponencial com \textit{jitter} e \textit{timeouts};
  \item Idempotência ponta a ponta, \textit{dedup} e \textit{replay} controlado;
  \item \textit{Bulkhead} e \textit{circuit breaker} para isolar falhas;
  \item Criptografia em trânsito/repouso, mascaramento de PII e controle de acesso por papéis;
  \item Observabilidade: métricas (latência, throughput, erro), logs estruturados e \textit{tracing} distribuído.
\end{itemize}

\textbf{Qualidades e cenários} (ATAM/SEI, ISO 42010/RM‑ODP):
\begin{itemize}
  \item Latência alvo por janela (diária/por sprint) e \textit{SLOs} de disponibilidade;
  \item Throughput sustentado/\textit{burst}, perda máxima tolerada e MTTR de reprocessamento;
  \item Portabilidade por conectores/contratos e neutralidade tecnológica.
\end{itemize}

\subsection{Caso 4 — Vistas de Observabilidade e Sinais para Docentes}

\textbf{Sinais derivados} (exemplos, não avaliativos):
\begin{itemize}
  \item PRs sem revisão \(>\) N horas; \textit{burst} de commits pré-review; divergência de \textit{branches};
  \item \textit{Issue aging} e WIP por equipe; tarefas bloqueadas por \(>\) N dias;
  \item Falhas de \textit{build/test} persistentes por janela; variação anômala de cadência de commits;
  \item Atualizações de documentação correlacionadas a marcos do TAPI.
\end{itemize}

\textbf{Materializações}:
\begin{itemize}
  \item Modelos de leitura por contexto (equipe, sprint, repositório);
  \item Encadeamentos explicáveis (link para eventos/artefatos subjacentes);
  \item Acesso por papéis (docente/coordenação) e trilhas de auditoria.
\end{itemize}

\textbf{Observações}:
\begin{itemize}
  \item Os sinais são \textit{insumos} para decisão docente e coordenação; não configuram avaliação discente;
  \item A rastreabilidade permite justificar cada sinal pelos eventos e chaves envolvidos.
\end{itemize}

\subsection{Resumo do Capítulo}

Apresentou-se uma arquitetura orientada a CEP organizada em quatro casos de modelagem: (i) esquema/contratos de eventos; (ii) rastreabilidade fim a fim; (iii) ingestão/CEP com táticas arquiteturais e qualidades; e (iv) vistas e sinais operacionais para apoio à decisão. Essa modelagem alinha-se à ISO/IEC/IEEE 42010 \cite{iso42010} e ao RM‑ODP \cite{iso10746}, priorizando observabilidade, rastreabilidade e governança de dados, sem qualquer mensuração de desempenho discente.
